% Options for packages loaded elsewhere
\PassOptionsToPackage{unicode}{hyperref}
\PassOptionsToPackage{hyphens}{url}
\PassOptionsToPackage{dvipsnames,svgnames,x11names}{xcolor}
\documentclass[
  11pt,
]{article}
\usepackage{xcolor}
\usepackage[margin=1in]{geometry}
\usepackage{amsmath,amssymb}
\setcounter{secnumdepth}{5}
\usepackage{iftex}
\ifPDFTeX
  \usepackage[T1]{fontenc}
  \usepackage[utf8]{inputenc}
  \usepackage{textcomp} % provide euro and other symbols
\else % if luatex or xetex
  \usepackage{unicode-math} % this also loads fontspec
  \defaultfontfeatures{Scale=MatchLowercase}
  \defaultfontfeatures[\rmfamily]{Ligatures=TeX,Scale=1}
\fi
\usepackage{lmodern}
\ifPDFTeX\else
  % xetex/luatex font selection
  \setmainfont[]{DejaVu Serif}
  \setmonofont[]{DejaVu Sans Mono}
\fi
% Use upquote if available, for straight quotes in verbatim environments
\IfFileExists{upquote.sty}{\usepackage{upquote}}{}
\IfFileExists{microtype.sty}{% use microtype if available
  \usepackage[]{microtype}
  \UseMicrotypeSet[protrusion]{basicmath} % disable protrusion for tt fonts
}{}
\makeatletter
\@ifundefined{KOMAClassName}{% if non-KOMA class
  \IfFileExists{parskip.sty}{%
    \usepackage{parskip}
  }{% else
    \setlength{\parindent}{0pt}
    \setlength{\parskip}{6pt plus 2pt minus 1pt}}
}{% if KOMA class
  \KOMAoptions{parskip=half}}
\makeatother
\usepackage{color}
\usepackage{fancyvrb}
\newcommand{\VerbBar}{|}
\newcommand{\VERB}{\Verb[commandchars=\\\{\}]}
\DefineVerbatimEnvironment{Highlighting}{Verbatim}{commandchars=\\\{\}}
% Add ',fontsize=\small' for more characters per line
\newenvironment{Shaded}{}{}
\newcommand{\AlertTok}[1]{\textcolor[rgb]{1.00,0.00,0.00}{\textbf{#1}}}
\newcommand{\AnnotationTok}[1]{\textcolor[rgb]{0.38,0.63,0.69}{\textbf{\textit{#1}}}}
\newcommand{\AttributeTok}[1]{\textcolor[rgb]{0.49,0.56,0.16}{#1}}
\newcommand{\BaseNTok}[1]{\textcolor[rgb]{0.25,0.63,0.44}{#1}}
\newcommand{\BuiltInTok}[1]{\textcolor[rgb]{0.00,0.50,0.00}{#1}}
\newcommand{\CharTok}[1]{\textcolor[rgb]{0.25,0.44,0.63}{#1}}
\newcommand{\CommentTok}[1]{\textcolor[rgb]{0.38,0.63,0.69}{\textit{#1}}}
\newcommand{\CommentVarTok}[1]{\textcolor[rgb]{0.38,0.63,0.69}{\textbf{\textit{#1}}}}
\newcommand{\ConstantTok}[1]{\textcolor[rgb]{0.53,0.00,0.00}{#1}}
\newcommand{\ControlFlowTok}[1]{\textcolor[rgb]{0.00,0.44,0.13}{\textbf{#1}}}
\newcommand{\DataTypeTok}[1]{\textcolor[rgb]{0.56,0.13,0.00}{#1}}
\newcommand{\DecValTok}[1]{\textcolor[rgb]{0.25,0.63,0.44}{#1}}
\newcommand{\DocumentationTok}[1]{\textcolor[rgb]{0.73,0.13,0.13}{\textit{#1}}}
\newcommand{\ErrorTok}[1]{\textcolor[rgb]{1.00,0.00,0.00}{\textbf{#1}}}
\newcommand{\ExtensionTok}[1]{#1}
\newcommand{\FloatTok}[1]{\textcolor[rgb]{0.25,0.63,0.44}{#1}}
\newcommand{\FunctionTok}[1]{\textcolor[rgb]{0.02,0.16,0.49}{#1}}
\newcommand{\ImportTok}[1]{\textcolor[rgb]{0.00,0.50,0.00}{\textbf{#1}}}
\newcommand{\InformationTok}[1]{\textcolor[rgb]{0.38,0.63,0.69}{\textbf{\textit{#1}}}}
\newcommand{\KeywordTok}[1]{\textcolor[rgb]{0.00,0.44,0.13}{\textbf{#1}}}
\newcommand{\NormalTok}[1]{#1}
\newcommand{\OperatorTok}[1]{\textcolor[rgb]{0.40,0.40,0.40}{#1}}
\newcommand{\OtherTok}[1]{\textcolor[rgb]{0.00,0.44,0.13}{#1}}
\newcommand{\PreprocessorTok}[1]{\textcolor[rgb]{0.74,0.48,0.00}{#1}}
\newcommand{\RegionMarkerTok}[1]{#1}
\newcommand{\SpecialCharTok}[1]{\textcolor[rgb]{0.25,0.44,0.63}{#1}}
\newcommand{\SpecialStringTok}[1]{\textcolor[rgb]{0.73,0.40,0.53}{#1}}
\newcommand{\StringTok}[1]{\textcolor[rgb]{0.25,0.44,0.63}{#1}}
\newcommand{\VariableTok}[1]{\textcolor[rgb]{0.10,0.09,0.49}{#1}}
\newcommand{\VerbatimStringTok}[1]{\textcolor[rgb]{0.25,0.44,0.63}{#1}}
\newcommand{\WarningTok}[1]{\textcolor[rgb]{0.38,0.63,0.69}{\textbf{\textit{#1}}}}
\usepackage{longtable,booktabs,array}
\usepackage{calc} % for calculating minipage widths
% Correct order of tables after \paragraph or \subparagraph
\usepackage{etoolbox}
\makeatletter
\patchcmd\longtable{\par}{\if@noskipsec\mbox{}\fi\par}{}{}
\makeatother
% Allow footnotes in longtable head/foot
\IfFileExists{footnotehyper.sty}{\usepackage{footnotehyper}}{\usepackage{footnote}}
\makesavenoteenv{longtable}
\setlength{\emergencystretch}{3em} % prevent overfull lines
\providecommand{\tightlist}{%
  \setlength{\itemsep}{0pt}\setlength{\parskip}{0pt}}
\usepackage{amsmath,amssymb,booktabs}
\setlength{\emergencystretch}{3em}
\usepackage{bookmark}
\IfFileExists{xurl.sty}{\usepackage{xurl}}{} % add URL line breaks if available
\urlstyle{same}
\hypersetup{
  pdftitle={Local Factor Rigidity and Mass--Defect Bounds for Langton's Ant},
  pdfauthor={Hao Ke},
  colorlinks=true,
  linkcolor={Maroon},
  filecolor={Maroon},
  citecolor={Blue},
  urlcolor={Blue},
  pdfcreator={LaTeX via pandoc}}

\title{Local Factor Rigidity and Mass--Defect Bounds for Langton's Ant}
\author{Hao Ke}
\date{}

\usepackage{fvextra}
\RecustomVerbatimEnvironment{verbatim}{Verbatim}{breaklines=true,breakanywhere=true,fontsize=\small}
\begin{document}
\maketitle
\begin{abstract}
Can the short turn patterns of a highway force a Langton ant to follow
that highway? We call a ten-turn window legal when it occurs in the standard
104-turn highway word, read cyclically. If every future window is legal and
the initial colouring has $m$ black cells, the ant reaches the standard
highway within $12m+212$ updates. The coefficient 12 is optimal.
For trajectories from any finite black set, the same language measures
departures from highway behaviour. Writing $B_t$ for the black cells at time $t$ and
$N_{10}(s,t)$ for the illegal windows wholly contained in $[s,t)$, we prove
\[
3(t-s)-26(|B_t|-|B_s|)\le351N_{10}(s,t)+322.
\]
Thus a trajectory with few illegal windows must create black cells at nearly
the highway rate. The additive constant 322 and the asymptotic mass-growth
slope are sharp. The attraction proof excludes arbitrarily repeated folds
in a finite graph of physically feasible words. The mass inequality uses a
separate finite-word potential and applies even to fragments with no infinite
legal continuation. We also compute the next illegal window and its complete
output state, describe the geometry of finite legal intervals, and show that
nine is the least window length giving a uniform legal-word mass-deficit bound.
\end{abstract}

\section{Introduction}\label{introduction}
The standard Langton highway repeats a sequence of 104 left and right turns.
A much shorter observation records only which ten-turn windows occur in that
cycle. Does matching this local vocabulary force the ant eventually to follow
the whole highway? The question involves the physical colouring as well as the
turn sequence: each visit flips a cell, so a word allowed by overlapping
windows may still request the wrong colour when it returns to an earlier site.

We answer this question for every finite initial black set. If all future
ten-turn windows belong to the highway vocabulary, the ant enters the standard
highway within $12m+212$ updates, where $m$ is the initial number of black
cells. The linear coefficient is optimal. On an infinite initial colouring,
two other eventual modes can occur. The proof identifies these possibilities
in a finite word graph and then excludes arbitrarily long combinations of
folds by explicit conflicts at repeated cell visits.

The same vocabulary yields a second result for trajectories with no promise
about their future turns. Every illegal ten-turn window contributes a fixed
allowance to a bound on the difference between elapsed time and black-cell
growth. Consequently a long interval with few violations must leave nearly
as many black cells as a highway would. This statement requires a separate
proof for finite words: a finite legal fragment need not admit a legal
infinite continuation. A finite graph potential supplies that proof.

After establishing these two results, we develop their algorithmic and
geometric consequences. We compute the first local violation and the complete
state it leaves, bound the geometry and black supply of finite legal
intervals, and construct examples that store and later read back black
memory. Nine-letter windows suffice for a uniform mass-deficit bound,
whereas ten are necessary to force eventual periodicity on arbitrary
colourings. These are different thresholds for different conclusions.

Gajardo's symbolic projection and forbidden-word framework {[}1{]}
provide the finite-word viewpoint. Gajardo and Mazoyer {[}4{]} relate
finite-state trace descriptions to restrictions on a head's return tours;
their one-dimensional equivalences are in Section 6 of the cited author text.
Gajardo, Lutfalla and Rao {[}2{]} study highway multiplicity for generalized
ant rules; Jillhewar's full-drift chronology {[}3{]} supplies the periodic
realization test used here. Our contribution is the concrete physical
classification of the short-factor language, its sharp quantitative bounds,
and the resulting local-event algorithms. The all-parameter exclusion of
repeated folds and the finite-word potential are the two main proof steps.

The paper first proves attraction and the mass inequalities. It then develops
exact event computation and finite-interval geometry. The appendices give the
explicit words, supply tables, and memory constructions used to attain or
illustrate the bounds.

\section{Model and principal results}\label{model-and-principal-results}

A state consists of the complete black set \(B\), a position in
\(\mathbb Z^2\), and a cardinal heading. White turns right (R), black
turns left (L); the cell then flips and the ant moves. Time denotes the
state before the update. Denote the state by $X=(B,p,h)$ and one update
by $F$. Headings N,E,S,W correspond to $h=0,1,2,3$ and vectors
$(0,1),(1,0),(0,-1),(-1,0)$. The update is invertible: step back along the
current heading, flip that cell, and undo the turn determined by its recovered
colour. Thus each complete state determines its turn sequence at all integer
times, including the nine turns preceding the input checkpoint.

\paragraph*{Turn language and convergence.}
Let \(W\) be the standard 104-letter turn word
in Appendix~\ref{exact-words} and let \(\mathcal L_k\) contain all its cyclic factors
of length \(k\). A \emph{factor} is a consecutive subword; cyclic factors
may cross the end of $W$ and continue at its beginning. A stream is
\(k\)-legal if every future length-\(k\) window belongs to
\(\mathcal L_k\). A physical stream is the turn sequence of an actual ant
on a specified colouring. The \emph{standard trace} repeats $W$, together
with its induced positions, headings and read colours, up to translation,
rotation and cyclic phase. The entry time $\tau(X)$ is the earliest update
from which this complete future trace is standard. The mass is $|B_t|$;
unvisited initial black cells remain part of that set. Let $H(X)$ denote
highway convergence, namely $\tau(X)<\infty$.

\paragraph*{Violation events and interval counts.}
The vocabulary $\mathcal L_{10}$ contains 84 distinct words. Forward legality
tests windows starting at times $t\ge0$. To define an event from the current
complete state alone, reconstruct its nine preceding turns by inverse dynamics.
Define $J_{10}(X)=1$ exactly when those nine turns followed by the current turn
form a word outside $\mathcal L_{10}$; otherwise $J_{10}(X)=0$.
This also defines the event at the input checkpoint without a supplied history.
The contained count $N_{10}(s,t)$ counts only forbidden windows wholly inside
$[s,t)$. It counts $J_{10}$ at ending times $s+9,\ldots,t-1$; it is zero when
$t-s<10$. For finite words, \emph{silent} means ten-legal; an event-free
trajectory interval excludes $J_{10}$ at each of its specified times.

\textbf{Theorem A (eventual modes).} Every physical ten-legal stream on
an arbitrary initial colouring eventually has least period 104, 28 or
24, with the respective words \(W,U,T\) printed in Appendix~\ref{exact-words}.

\textbf{Theorem B (finite-mass attraction).} For a complete finite
initial black set of cardinality \(m\), a ten-legal future reaches the
standard trace by \[\tau(X)\le 12m+212.\] The coefficient 12 is optimal:
a family with \(m_n=2n+13\) has \(\tau_n=24n-3\) for every \(n\ge1\).
The intercept 212 is a proved upper bound.

\textbf{Theorem C (arbitrary-colouring threshold).} An explicit
computable infinite colouring has a nine-legal stream that is not
eventually periodic. Thus ten is the least factor length forcing
eventual periodicity on arbitrary colourings. This sharpness statement
uses infinite support.

\textbf{Theorem D (local event).} Let \(J_{10}\) test the true preceding
nine turns and the current turn for membership outside
\(\mathcal L_{10}\), recovering the past by inverse dynamics. On every
finite input, existence of a future \(J_{10}\) is decidable in
polynomial bit time for sparse binary coordinates.

\textbf{Theorem E (mass--defect balance).} On every interval \([s,t)\)
of every finite-support trajectory,
\[3(t-s)-26\bigl(|B_t|-|B_s|\bigr)\le351N_{10}(s,t)+322,\] where
\(N_{10}\) counts forbidden ten-letter windows wholly contained in the
interval. The constant 322 is attained by a 90-step witness; the mass
slope \(3/26\) is asymptotically sharp. The penalty 351 is not asserted
optimal.

Theorems A--C concern infinite legal futures; Theorem D uses the
finite-mass attraction bound and a complete-field periodic-tail test. Theorem E instead rests on a complete finite-word
potential certificate. This separates their logical dependencies while
using the same local language.

\section{From local language to eventual
attraction}\label{from-local-language-to-eventual-attraction}

\subsection{Finite-word feasibility and graph coverage}\label{ten-2}

For any finite prescribed turn word, the geometric path is fixed after
choosing an initial pose. It is physically realizable if and only if
consecutive reads of each same physical cell have alternating letters.
Necessity is the flip rule. For sufficiency, make a visited cell
initially black exactly when its first requested letter is L; leave
every unvisited cell white. Induction over the finite word proves
execution. Gajardo {[}1{]}, Section 3, pp.\ 60--62, gives this
first-read assignment and its forbidden-word formulation.

Let \(E\) contain every physically feasible 150-letter word whose
10-factors lie in \(\mathcal L_{10}\). Construct a graph \(G\) with edge
\(v_0\cdots v_{148}\to v_1\cdots v_{149}\) for each \(v\in E\). The length 150 provides a finite cover containing the candidate periods
104, 28 and 24 within a window. Its role is to expose a graph whose cyclic
components and physical exclusions can be completely certified; no minimality
of this window length is needed.
Complete enumeration follows all whitelist-compatible extensions and
rejects exactly the same-cell alternation conflicts. Induction on word
length proves completeness of that enumeration.

Any actual infinite 10-legal stream gives an infinite path in \(G\),
since each of its length-150 windows is feasible under its own actual
execution. The graph can admit paths that fail longer physical
consistency. It is an overapproximation: every physical stream is represented, and
additional paths are eliminated by the physical conflicts below.

Complete enumeration gives 6,605 edge words and 6,501 vertices. The
finite certificate, identified in Section~\ref{reproducibility-and-conclusion},
checks that the graph has precisely three cyclic strongly connected components:

\begin{longtable}[]{@{}lrrl@{}}
\toprule\noalign{}
Component & Vertices & Internal edges & Structure \\
\midrule\noalign{}
\endhead
\bottomrule\noalign{}
\endlastfoot
\(\mathcal W\) & 104 & 104 & exact standard simple cycle \\
\(\mathcal T\) & 24 & 24 & exact \(T\) simple cycle \\
\(\mathcal H\) & 366 & 370 & not a simple cycle \\
\end{longtable}

There are no other cyclic components, including self-loops. The first
two component edge sets equal all cyclic length-150 factors of \(W\) and
\(T\). Strongly connected components and reachability are checked on the
full graph, not only on a pruned experimental sample.

The non-simple component is crucial: ordinary finite-state periodicity
reasoning cannot dispose of it. It is dealt with next by an all-length
physical conflict, rather than by assuming its paths are periodic.

\subsection{Exact two-loop description of the non-simple
component}\label{ten-3}

An explicitly tabulated map assigns each of the 366 vertices of
\(\mathcal H\) a state in \(\{0,\ldots,101\}\). It is a
label-preserving homomorphism onto this explicit graph:

\begin{itemize}
\tightlist
\item
  for each \(0\le i<102\), the edge labelled \(B_i\) leads from \(i\) to
  \(i+1\) modulo 102;
\item
  one additional edge, labelled L, leads from 0 to 75.
\end{itemize}

Here \(B\) is the 102-letter word in the appendix, \(B_0=\mathrm R\),
and the alternative short loop is \(A=\mathrm L B_{75}\cdots B_{101}\),
of length 28. Every one of the 370 internal edges is checked against
this formula. The homomorphism is used only in the forward direction:
every infinite internal path projects to a word that, after at most 101
letters, is an arbitrary concatenation of \(A\) and \(B\).

The edge label is the letter appended to the 149-letter history. It is
shifted 149 places from the original stream. This finite shift is
harmless for eventual-mode classification. Subsection \ref{ten-5}
instead uses the current letter \(v_0\) of a future-window vertex,
149 positions before this appended-edge label.

The exact geometric data, starting north at the origin, are

\[A:(0,-4),\ h=0;\qquad B:(1,11),\ h=2.\]

Thus \(A\) translates and \(B\) reverses the heading. The \(A\) word is
a cyclic phase of \(U\).

\subsection{Three foldbacks are impossible, at every pair of
lengths}\label{ten-4}

\textbf{Lemma 1.} For every pair of integers \(n,m\ge0\), the finite
word \(B A^n B A^m B\) is physically infeasible, independently of its
initial colouring.

In each case below there are two consecutive visits to a cell \(q\) both
requesting R:

\begin{longtable}[]{@{}llrr@{}}
\toprule\noalign{}
Parameter range & \(q\) & Earlier time & Later time \\
\midrule\noalign{}
\endhead
\bottomrule\noalign{}
\endlastfoot
\(n=0\) & \((0,2)\) & 32 & 184 \\
\(n\ge1,\ m\ge n\) & \((0,2)\) & 32 & \(184+56n\) \\
\(n\ge1,\ 0\le m<n\), \(r=n-m\) & \((1,4r)\) & \(25+28r\) &
\(205+28(n+m)\) \\
\end{longtable}

All times are pre-update indices in the displayed finite word. It
suffices to verify their consecutiveness, not that the rest of the
candidate word is feasible.

\textbf{Proof.} Let the five block boundaries start at times 0, 102,
\(102+28n\), \(204+28n\), and \(204+28(n+m)\). Their respective poses
are \((0,0),N\); \((1,11),S\); \((1,11+4n),S\); \((0,4n),N\); and
\((0,4(n-m)),N\).

Local word coordinates below are based at the origin heading north. A word
placed at $p$ heading north sends a local point $q$ to $p+q$; heading south
sends it to $p-q$. In particular, copy $j$ of the first $A$ run starts at
$(1,11+4j),S$, so its local column $x=0$ is global column $x=1$, and local
height $y$ becomes $11+4j-y$. The second $A$ run instead starts north and
uses the unrotated local axes.

The following finite coordinate facts are computed directly from the two
displayed words and checked independently: \(A\) has \(-3\le x\le0\); on
its column \(x=0\) the pairs (index, height, letter) are

\[ (0,0,L),(7,-1,R),(8,-2,R),(15,-3,L),(16,-2,L),(27,-5,L).\]

In \(B\), the only visit to \((0,2)\) is index 32, R. All \(B\) visits
at \(x=0\) have \(y\le2\). Visits to \((1,9)\) are indices 82,R and
90,L; visits to \((1,4)\) are 45,L and 53,R; visits to \((1,8)\) are
73,L and 81,R. On \(x=1\) its maximum height is 12.

If \(n=0\), the reflected second \(B\) reads \((0,2)\) at its local
\((1,9)\), index 82. The previous visit is the first \(B\)'s index 32.
This gives the first row.

If \(n\ge1\), the first \(A\) run lies at \(x\ge1\), so cannot revisit
\((0,2)\). In the second \(B\), a visit there would require local
\((1,9+4n)\), above every \(x=1\) visit. In the second \(A\) run, the
first matching copy is number \(n-1\) (zero-based), at local index 8,
with letter R. All earlier copies miss: on \(x=0\) their required local
height is below -5. The visit time is \(204+28n+28(n-1)+8=184+56n\).
This copy exists when \(m\ge n\).

Finally let \(r=n-m>0\). For \(r=1\) or 2, the last earlier visit to
\((1,4r)\) is the first \(B\)'s index 53 or 81, both R. The first \(A\)
run lies at heights at least 11 and cannot reach those two cells. For
\(r\ge3\), the last visit in the first \(A\) run is copy \(r-3\), local
index 7, R, at time \(102+28(r-3)+7=25+28r\). The column-\(x=0\) list
for \(A\) shows no later visit to that cell in the run; any visit in the
first \(B\) precedes it. A visit by the second \(B\) would require local
\((0,11+4m)\), impossible since its column-\(x=0\) maximum is 2. The
second \(A\) run lies at \(x\le0\). The third \(B\), based at
\((0,4r)\), first reaches \((1,4r)\) at index 1, R. This proves the
third row and the lemma. \(\square\)

An infinite concatenation of \(A\) and \(B\) with infinitely many
\(B\)'s contains a forbidden three-\(B\) subword. Hence it has only
finitely many \(B\)'s and is ultimately \(A^\omega\). This argument
disposes of arbitrary interleavings and unbounded run lengths, not only
tested values of \(n,m\).

An infinite path in a finite graph eventually remains in one cyclic
strongly connected component. In \(\mathcal W\) and \(\mathcal T\) it is
forced periodic by the simple cycles. In \(\mathcal H\) the preceding
homomorphism and Lemma 1 force an \(A\)-periodic suffix. This proves
Theorem A.

\subsection{Finite support and the transient bound}\label{ten-5}

For a heading-resetting periodic word of length \(P\) with nonzero drift
\(d\), partition its phase positions by differences in \(\mathbb Z d\).
For a class write \(q_i=b+a_i d\). Visits to \(b+nd\) are eligible when
\(a_i\le n\), and their chronological order is decreasing \(a_i\), then
increasing \(i\). A cell's required initial colour is the first letter
of that ordered eligible list. This is the full-drift chronology method
of Jillhewar (2026), Theorem 3.1 and Remark 3.2.

The words \(U\) and \(T\) pass all per-cell alternation tests but have
respectively five and two stabilized classes starting L. Each such class
demands infinitely many DISTINCT initially black cells. Neither word can
be executed forever from finite support. Any finite prelude preserves
finite support, so these modes also cannot occur as eventual tails of a
finite-support orbit. Consequently every finite-support infinite
10-legal path eventually enters \(\mathcal W\).

The complete graph gives 925 vertices that can reach \(\mathcal W\).
They exclude \(\mathcal H\) entirely and include \(\mathcal T\). Apart
from \(\mathcal W\) and \(\mathcal T\), they are acyclic. A path cannot
leave a strongly connected component and later return to it. The current
letters of a run of \(L\) successive vertices in \(\mathcal T\) form a
length-\(L\) prefix of a phase of \(T^\omega\).

Let \(b_j(L)\) count distinct cells whose first read is L in this prefix
starting at phase \(j\). If the complete black mass at the run's start
is \(m_T\), physical execution requires \(b_j(L)\le m_T\). The count
includes only first reads relative to that run; cells written during the
run cannot pay for them.

\textbf{Lemma 2.} For all phases \(j\) and all integers \(L\ge0\),

\[L-12b_j(L)\le5.\]

\textbf{Proof.} Chronological classes give the first-L read times as a
finite set \(F_j\) and two arithmetic rays \(r_{j,1}+24\mathbb N\),
\(r_{j,2}+24\mathbb N\). Their exact table is in the appendix. The sets
count distinct physical cells, not repeated L reads. Put
\(D_j=\max(\{0\}\cup\{f+1:f\in F_j\}\cup\{r_{j,1},r_{j,2}\})\). For
every \(L\ge D_j\), directly from these rays,

\[b_j(L+24)=b_j(L)+2.\]

Thus \(L-12b_j(L)\) is 24-periodic for all \(L\ge D_j\). Checking the
finite interval \(0\le L<D_j+24\) covers all lengths. Both
implementations derive the class table independently and verify that the
maximum across all phases is 5. This is an exact finite reduction with
its all-length periodicity proof, not a finite-prefix extrapolation.
\(\square\)

Before its first \(\mathcal T\) vertex a path uses only the 54 acyclic
strict ancestors of \(\mathcal T\). Give such a vertex \(v\) weight

\[c(v)=1+12\delta(v_0),\qquad \delta(R)=1,\ \delta(L)=-1.\]

Complete acyclic dynamic programming gives maximum total weight 70 for a
path ending at \(\mathcal T\). Therefore if \(t_{pre}\) updates precede
that run and their net black change is \(\Delta\), then
\(t_{pre}+12\Delta\le70\) and \(m_T=m+\Delta\). After departure from
\(\mathcal T\), the maximum acyclic path length before \(\mathcal W\) is
137. The exit-producing update from a \(\mathcal T\) vertex is included
in its \(T\)-prefix, not omitted. By Lemma 2,

\[\tau\le t_{pre}+(12m_T+5)+137\le12m+70+5+137=12m+212.\]

Paths avoiding \(\mathcal T\) altogether have length at most 167 before
\(\mathcal W\), also below this bound. A path whose initial vertex is
already in \(\mathcal W\) has zero such entrance time. Once the path
first reaches \(\mathcal W\), leaving it would preclude a later return;
because its eventual component must be \(\mathcal W\), it never leaves.
Its future is standard from that time.

These are future-window vertices at times starting at zero, so the transient bound needs no additional
window warm-up. Every
complete initial black cell remains in the physical state; only its
count is used in this inequality. Theorem B's upper bound follows.

\subsection{An explicit family attaining coefficient
twelve}\label{ten-6}

Define

\[E=\{(-2,-1),(-1,0)\},\]

\[\begin{split}D=\{&(-1,-4),(-1,-3),(0,-5),(0,-3),(0,-2),(1,-6),\\
 &(1,-5),(1,-4),(1,-3),(2,-3),(2,-2)\}.\end{split}\]

\[R=\{(-2,0),(-1,-1),(-1,0),(0,-1),(1,-1),(1,0)\}.\]

For \(n\ge0\), put

\[B_n=E\ \cup\ \{(-2,-2j),(0,-2j-1):1\le j\le n\}\ \cup\ (D+(0,-2n)).\]

These sets are disjoint. Start north at the origin with exactly \(B_n\)
black and all other cells white. Then \(|B_n|=2n+13\).

The word \(T\) has read footprint within \([-2,1]\times[-3,0]\), net
displacement \((0,-2)\), and resets heading. For every \(n\ge1\) its
required initial black cells in that footprint are precisely
\(E\cup\{(-2,-2),(0,-3)\}\). All later rail cells and all cap cells are
outside the footprint. The complete 24-step parity of visits, followed
by set cancellation, gives the identity

\[F^{24}(B_n,0,N)=\big((B_{n-1}+(0,-2))\mathbin{\dot\cup}R,\ (0,-2),N\big).\]

This is an identity of FULL states, not a core-only recurrence. A direct
all-parameter derivation is: translate the rail terms with
\(j=1,\ldots,n-1\) by -2; they become the original terms
\(j=2,\ldots,n\). The cap becomes \(D+(0,-2n)\). Only the fixed finite
set \(E\), the first rail pair, and the translated \(E\) remain to
compare, and their difference with \(R\) equals the exact net toggle set
of \(T\). This finite equality is checked. It also covers \(n=1\) where
the shorter rail is empty.

The residual \(R\) is at heights -1 and 0. Every future repeated \(T\)
segment based at height -2 or less remains strictly below it. The
\(n=0\) complete state is the clean seed of the standard phase \(Z\) in
the appendix; its full infinite execution is verified by chronology and
an independent epoch-breakpoint formula. Its phase footprint has maximum
height zero and drift \((-2,-2)\), so its continuation from height -2
also stays below \(R\). Induction therefore retains and guards all
successive residuals and proves the exact infinite word

\[T^n Z^\omega.\]

The factors of \(T^2Z^2\) verify every possible length-10 junction;
\(T^\omega\) and \(Z^\omega\) are themselves 10-legal. The last three
letters of \(T\) and \(Z\) agree, but the last four do not. Since \(Z\)
has least period 104, any earlier standard future agreeing with
\(Z^\omega\) must be its unique phase-aligned left extension. Hence the
first standard time is

\[\tau_n=\max(0,24n-3).\]

For \(n\ge1\), \(\tau_n=12|B_n|-159\). An affine upper bound with slope
less than 12 and fixed intercept fails for large \(n\). This proves the
sharp coefficient, without claiming 212 is sharp. Every member of this
family eventually reaches the standard highway.

\subsection{Nine factors admit an aperiodic infinite-background
spiral}\label{ten-7}

Let \(U\) and \(C\) be the 28- and 89-letter words in the appendix and
let \(n_i=5\cdot2^i\). Consider

\[a=\prod_{i=0}^{\infty} U^{n_i}C.\]

The prescribed geometry is explicitly computable, starting at the origin
heading north. Its all-future physical feasibility follows from the next
separation proof; it is not assumed.

A \(U\) block restores heading, moves \((0,-4)\), and has footprint in
\([0,3]\times[-5,0]\). A \(C\) block rotates the heading
counterclockwise, moves \((13,-4)\), and has footprint in
\([0,14]\times[-4,0]\). The finite word \(UCU\) is feasible. Every
length-9 factor of \(U^2CU^2\) is in \(\mathcal L_9\).

The start heading of macro \(i\) is \(-i\) modulo four. For \(S=16^k\),
the four consecutive start positions are

\[p_{4k}=(-8S+8,4S-4),\quad p_{4k+1}=(-8S+21,-16S-8),\]
\[p_{4k+2}=(32S+25,-16S+5),\quad p_{4k+3}=(32S+12,64S+9).\]

These follow by the two block displacements; one group of four returns
the heading and replaces \(S\) by \(16S\). Entire macro read footprints
lie in the respective rectangles

\begin{longtable}[]{@{}lll@{}}
\toprule\noalign{}
Index & x interval & y interval \\
\midrule\noalign{}
\endhead
\bottomrule\noalign{}
\endlastfoot
\(4k\) & \([-8S+8,-8S+22]\) & \([-16S-8,4S-4]\) \\
\(4k+1\) & \([-8S+21,32S+25]\) & \([-16S-8,-16S+6]\) \\
\(4k+2\) & \([32S+11,32S+25]\) & \([-16S+5,64S+9]\) \\
\(4k+3\) & \([-128S+8,32S+12]\) & \([64S-5,64S+9]\) \\
\end{longtable}

Each rectangle is disjoint from all nonadjacent earlier rectangles. For
the first row, \(k\ge1\), earlier rectangles through \(4k-2\) have
\(x\ge-S/2+8\), while the new maximum is \(-8S+22\); use \(15S>28\). For
the second row, earlier rectangles through \(4k-1\) have \(y\ge-S-8\),
while the new maximum is \(-16S+6\); use \(15S>14\). For the third row,
earlier rectangles through \(4k\) have \(x\le2S+25\), while the new
minimum is \(32S+11\); use \(30S>14\). For the fourth row, earlier
rectangles through \(4k+1\) have \(y\le4S+9\), while the new minimum is
\(64S-5\); use \(60S>14\). The first two empty-earlier-set base cases
are automatic. These inequalities hold for every \(S=16^k\) in the
stated ranges.

It remains to control adjacent macros, which are not claimed disjoint.
In local coordinates of the earlier macro, the next macro begins at
\((13,-4n_i-4)\) heading west. Its first \(U\) block can overlap the
earlier \(C\), but every later \(U\) block has \(x\ge17\) while \(C\)
has \(x\le14\); the next \(C\) is still farther right. The preceding
\(U\) rail has \(x\le3\), so is disjoint from the next macro. Within one
macro, \(C\) can meet only the LAST \(U\) block: every earlier \(U\)
block has minimum height at least \(-4n_i+3\), above \(C\)'s maximum
\(-4n_i\). Thus all interactions with \(C\) are contained in the
verified word \(UCU\). Other repetitions are portions of the physically
feasible \(U^\omega\) trace. This proves alternating letters at every
consecutive same-cell pair in the entire infinite word.

Colour a cell initially black precisely when its first prescribed read
has letter L; leave unvisited cells white. The proved consistency makes
this a genuine infinite classical-ant execution by induction over time.
This colouring is computable: for a queried cell of infinity-norm \(M\),
choose \(k\) with \(16^k>M+100\). Every rectangle from macro \(4k\)
onward lies outside \([-M,M]^2\) in at least one coordinate, so only the
first \(4k\) finite macros need be searched. It has infinite black
support, since interior portions of arbitrarily long \(U\) rails demand
fresh initial black cells.

The word \(D_{bad}=\mathtt{RLRRRRLRRL}\) is absent from
\(\mathcal L_{10}\) and occurs exactly once in each connector \(C\), at
offset 11. The finite check of \(U^2CU^2\) excludes other occurrences at
either junction, and \(U^\omega\) has none. Successive occurrences are
separated by \(89+28n_{i+1}\), which grows without bound. An eventually
periodic sequence with infinitely many occurrences of a fixed finite
factor has bounded gaps between them. Hence \(a\) is not eventually
periodic. This proves Theorem C, including its explicitly
infinite-background boundary.

\subsection{A radius-nine target decidable without waiting for a remote
collision}\label{ten-8}

The inverse of a literal step is unique: from post-pose \((p,h)\) set
\(q=p-\mathrm{dir}(h)\), undo the flip at \(q\), and undo the turn
dictated by that recovered colour. Nine inverse steps read only cells
within distance nine of the current head. Therefore the
preceding-nine-turn test \(J_{10}\) is an intrinsic finite-local
predicate, including at input time zero.

Let \(M=12|B_0|+212\). The following total algorithm acts on an
arbitrary finite input, with no legality promise:

\begin{enumerate}
\def\labelenumi{\arabic{enumi}.}
\tightlist
\item
  Recover the true nine-turn past. Simulate up to time \(M\), checking
  \(J_{10}\) before each update. Return YES with its first time if
  encountered.
\item
  At time \(M\), test the complete-state predicate \(P_{now}\) for
  permanent standard P104. If it fails, return YES (a future event is
  forced by Theorem B). This branch decides existence without simulating
  to the event; the exact first-time algorithm is developed in Section~\ref{macro-5}.
\item
  If \(P_{now}\) holds, simulate nine additional updates to check all
  windows crossing this boundary. Return YES if one is bad; otherwise
  return NO. All subsequent windows are inside the certified standard
  future.
\end{enumerate}

Why is the second step sound? Absence of every future \(J_{10}\) implies
legality of every forward length-10 window from time zero. Theorem B
would force \(P_{now}\) by time \(M\), contradicting the failed test. A
far-away blocker cannot be omitted from that test.

The exact \(P_{now}\) test first simulates 104 turns on a copy. If they
do not equal any phase of \(W\), it returns false. Otherwise let \(K\)
be the unique clean finite seed for that placed phase, computed by the
chronological class formula. Let \(q_i+\mathbb N d\) be its 104 future
phase rays. The actual complete field realizes that future if and only
if

\[(B\triangle K)\cap\bigcup_{i=0}^{103}(q_i+\mathbb N d)=\varnothing.\]

Sufficiency follows by first-difference induction; necessity follows at
the first read of an initial differing cell. Testing finitely many
complete-field points against 104 integer rays is exact, including
arbitrarily distant markers. No repeated-prefix heuristic is used.

During the \(O(m)\) ordinary updates the complete field has \(O(m)\)
cells, and coordinate bit lengths grow by only \(O(\log(m+1))\) beyond
their input sizes. Exact finite-set operations and integer ray tests
therefore give a deterministic polynomial bit-time implementation (a
balanced search tree can be used instead of hashing). The implementation
uses Python integers and finite sets; even worst-case linear set lookups
remain polynomial. There is no claim of a polynomial bound for the
numerical first-event time or of a bound polynomial in the bit length of
a COMPRESSED black count.

\subsubsection{A fixed-mass late-event
control}\label{a-fixed-mass-late-event-control}

\texttt{verification/ten/certificates/far\_marker.json} gives a complete
27-black standard state at a mature phase. Add one black marker at
\((2n,2n)\) for any \(n\ge20\). The new complete mass is always 28. The
first contact is exactly at update \(104n-92\), and it changes a
prescribed R into L, forming the forbidden ten-factor
\texttt{RRRLLLLRLL}.

The all-parameter first-contact calculation is: for the supplied phase
positions \(q_i\), write \(q_i=s_i(2,2)\) whenever possible. The first
visit to the marker has time \(104n+\min_i(i-104s_i)\). The minimum is
-92, attained at phase 12, and all relevant epochs are nonnegative for
\(n\ge20\). The marker is outside the finite true nine-step inverse
history and the seed. Until contact, first-difference coupling preserves
the entire standard path. The standard cyclic preceding nine letters and
the changed L give exactly the stated bad factor.

For \(n=10^{12}\) the first event is at \(103999999999908\), but the new
decision algorithm returns a forced-future YES after only 548
main-trajectory updates. The extra operations are nine inverse updates
and the constant lookahead and ray tests, not simulation to the event.
This family illustrates the distinction between reachability decision
and event-time observation. Its coordinates encode the long waiting time.

\subsection{Finitely many violations and standard
convergence}\label{finitely-many-violations-and-standard-convergence}

For finite inputs, a standard future has no intrinsic violations after
its first nine updates. Conversely, after the last violation every
forward window is legal; Theorem B applied to that complete state gives
a standard future. Hence \[H(X)\iff \#\{t\ge0:J_{10}(F^tX)\}<\infty.\]
Theorem D decides each future-event query. This equivalence identifies
the additional global assertion needed for convergence: finiteness of
the event chain, rather than decidability of an individual query.

\section{Mass and defects on arbitrary
intervals}\label{mass-and-defects-on-arbitrary-intervals}

For a finite word \(v\), put \(g(v)=\#R-\#L\) and \(D(v)=3|v|-26g(v)\).
\textbf{Theorem (finite-word deficit).} Every physically feasible,
ten-legal finite word satisfies \(D(v)\le322\), with no
infinite-extension premise. We derive the all-input inequality of Theorem E from this finite-word
bound.

\subsection{Exact finite-word feasibility and enumeration
coverage}\label{mass-2}

A prescribed turn word fixes its geometric path from $(0,0)$, heading N. It is
physically feasible exactly when successive reads of each physical cell
request opposite turns. Necessity is the flip rule. For sufficiency, put
an initial black cell precisely at each visited position whose first
request is L, and leave all other positions white. Induction realizes
the entire word. The required seed is finite. Arbitrary poses follow by
translation and quarter turns.

This elementary criterion is prior methodology, as in Gajardo's symbolic
forbidden-word treatment. It implies in
particular that finite-word feasibility is the same whether initial
colourings are required to be finite or are arbitrary.

Starting at the empty word, branch on BOTH turns. Reject an extension
exactly if its new ten-factor is forbidden, or its current cell was
previously read with the same letter. Induction on length proves
complete enumeration of physically feasible ten-legal words through
depth 300. Words of lengths below ten are initially unrestricted by
the standard factor table.

The enumeration has no seed-size or spatial cutoff. Its certificate
checks every accepted word and the integer potentials on the complete graph;
Section~\ref{reproducibility-and-conclusion} describes the two reconstructions.

\subsection{The finite graph certificate covers arbitrarily long finite
words}\label{mass-3}

Form a graph from every accepted word of length 300: its 299-prefix is
the source, its 299-suffix the target, and the appended letter labels
the edge. There are exactly 29,334 vertices and 29,559 edges. Assign the
weight

\[
w(R)=-23,\qquad w(L)=29,
\]

so that $D(v)$ is the sum of its letter weights.

The certificate assigns a nonnegative integer $h(u)$ to each vertex, with
maximum 251, and verifies every edge inequality

\begin{equation}\label{eq:mass2}
h(u)\ge w(c)+h(v)\quad\text{for every edge }u\xrightarrow{c}v
\end{equation}
and every initial-vertex inequality

\begin{equation}\label{eq:mass3}
D(u)+h(u)\le284.
\end{equation}

These are direct integer comparisons; an optimizer's output is not
trusted. The potential generators are only proposal generators. Their
resource caps raise an error with no mathematical conclusion. An
independently supplied damaged potential is rejected without repair.

For any physically feasible ten-legal finite word $z$ of length
$n\ge300$, all its 300-windows belong to this graph. They define
a path starting with its first 299 letters $u_0$. Summing \eqref{eq:mass2} over the
appended letters gives

\[
D(z)\le D(u_0)+h(u_0)-h(u_{\rm final})\le284,
\]

using \eqref{eq:mass3} and $h\ge0$. This is an overapproximation
argument: not every graph path need be physically realizable, but every
physical word is covered. No extension to an infinite ten-legal path is
assumed. There is no misuse of eventual-cycle classification for a
finite segment.

For all lengths 0 through 299, exhaustive enumeration gives maximum 322.
Combining the two finite kernels proves the finite-word deficit theorem for all finite
lengths.

The source witness fixes SHA256 hashes for the sorted depth-300 words
and the potentials in lexicographic vertex order.

The finite sets are reconstructed deterministically and checked against
these hashes and the local inequalities. The mass certificate also supplies
the following equality witness and its complete minimal seed.

\subsection{Sharpness}\label{mass-4}

The word

\begin{verbatim}
LLLLRLRRRRLRLLRRLLLLRRLLRRRRLLRRLRLLRLLRLRLLRRRRLLRLRRRRLRLLRRLLLLRRLLRRRRLLRRLRLLRLLRLRLL
\end{verbatim}

has length 90, 44 right turns and 46 left turns. It is physically
feasible and all of its 81 ten-factors are allowed. Its unique minimal
first-read seed has 28 black cells, obtained by blackening precisely the
first-L sites of the displayed word. Literal execution leaves 26 black cells.
Thus

\[
3\cdot90-26(26-28)=322,
\]

proving the additive constant in the finite-word deficit theorem and in Theorem E, with wholly
contained windows, cannot be decreased.

The standard word has 104 updates and 12 more R than L. Its certified
finite-support realization supplies physically feasible words $W^k$ for
every $k\ge1$. Thus any uniform lower estimate
$g\ge\rho n-C$ with $\rho>12/104=3/26$ fails as $k$ tends
to infinity. Permanent repetition of the standard highway supplies the
universal $k$ quantifier; its net mass gain is exactly 12 per block.

\subsection{Derivation of the unrestricted inequality}\label{mass-5}

In any actual finite orbit segment, every right turn increases complete
black mass by one and every left turn decreases it by one. Hence

\begin{equation}\label{eq:mass4}
g(a_s\cdots a_{t-1})=|B_t|-|B_s|.
\end{equation}

Remove each letter that ends a forbidden ten-window wholly contained in
the segment. These $N=N_{10}(s,t)$ single-letter positions separate at most
$N+1$ consecutive remaining runs. Every ten-window wholly inside a
remaining run was wholly inside the original segment and cannot be
forbidden: otherwise its last position would have been removed. Each run
is a genuine physical subword and is therefore covered by the finite-word deficit theorem. An
empty run has deficit zero and may harmlessly be charged the same upper
bound.

Each removed letter has deficit at most 29. Additivity of $D$, the finite-word deficit theorem
and \eqref{eq:mass4} yield

\[
D(\text{total})\le322(N+1)+29N=351N+322,
\]
proving Theorem E. This does not require disjoint forbidden windows.
Removing their ending letters handles overlaps exactly.

\subsection{Density consequences}\label{density-consequences}

Taking \(s=0\) and dividing by \(t\) gives two immediate consequences:
if \(|B_t|=o(t)\), then \(\liminf N_{10}(0,t)/t\ge1/117\); if
\(N_{10}(0,t)=o(t)\), then \(\liminf |B_t|/t\ge3/26\). These quantify
the cost of maintaining low mass or a nearly legal stream.

\section{Nine factors suffice for mass balance}\label{sec:nine-mass}
The language threshold for a uniform mass bound differs from the threshold
for eventual periodicity. Put
$D(v)=3|v|-26(\#R(v)-\#L(v))$. Every finite physical nine-legal word satisfies
\[
 D(v)\le1004.
\]
Consequently every complete finite-support orbit interval satisfies
\[
 3(t-s)-26(|B_t|-|B_s|)\le1033N_9(s,t)+1004,
\]
where $N_9$ counts forbidden nine-windows wholly contained in the interval.
The constant 1004 is not claimed optimal. The stronger ten-language
constant 322 remains valid on its narrower class.

\subsection{Finite certificate and all-length coverage}
Two complete enumerators produce all 164,283 feasible nine-legal words of
length 150. They also cover all shorter words, whose maximum deficit is
322, including words with no legal length-150 extension. The 149-overlap
graph has 157,114 vertices. Every longer real word is its initial history
followed by a graph path; the converse need not hold.

There are three cyclic components: simple cycles of lengths 104 and 24,
and a component with 34,712 vertices. A label-preserving graph map sends
the latter to a 1,213-state quotient: every source edge maps to an edge with
the same R/L label. This property preserves all physical paths even though
the quotient may admit additional paths. The refinement procedure finds the
map; checking its complete edge table establishes the property used here. Fourteen explicitly checked forbidden words each demand two
consecutive same-colour reads of one site and are physically impossible;
Appendix~\ref{app:nine-forbidden} prints the words and conflict indices.
The reachable product avoiding these words has 9,154 states and 9,550 edges.
Every quotient state is allowed as an entry with empty suffix history,
which preserves coverage even when a conflict crosses component boundaries.

An integer potential $0\le y_u\le306$ satisfies
\[
 y_u-y_v\ge\omega(c)\quad(u\mathrel{\mathop{\longrightarrow}^{c}}v),
 \qquad \omega(R)=-23,\quad\omega(L)=29.
\]
All inequalities are checked after reconstructing the potential.
Telescoping bounds the deficit of every internal path by 306.
The two simple-cycle total weights are 0 and $-32$; their maximal
subpath weights are 237 and 94. Removing complete laps proves these
bounds for arbitrarily many repetitions.

Give each component its internal budget and each entry its maximum
149-letter history deficit. A topological dynamic program on the complete
condensation graph adds crossing-edge weights and the next internal budget.
Every real path visits each component at most once. Checking all initial
and cross-edge inequalities gives the total bound 1004, including both
ends of the word. This establishes the all-length claim without extending
a finite simulation cutoff. Finally, remove the last letter of every bad
nine-window. The at most $N_9+1$ remaining consecutive runs are nine-legal;
each removed letter costs at most 29. Additivity yields
$1004(N_9+1)+29N_9$ and proves the unrestricted inequality.

\subsection{The sharp window threshold uses finite inputs}
Let
\[
 u=\mathrm{LLLLRRLLRRRRLLRRLRLLRR},\qquad d=(-3,1),
\]
and define
\[
 \begin{aligned}
 A&=\{(0,-1),(1,-1),(2,-1)\},\\
 H&=\{(-1,1),(-2,1),(0,0),(-1,0),(-1,-1)\}.
 \end{aligned}
\]
All cyclic eight-factors of $u$ are standard, and it contains eleven
letters of each colour. For every $n\ge0$ the complete finite source
\[
 B^{(n)}=A\cup\bigcup_{a\in H}\{a+jd:0\le j<n\}
\]
has $3+5n$ black sites and, from the origin facing north, executes $u^n$.
At update $22n$ it first departs from the proposed repetition.

For the all-parameter proof, partition phase positions into full
translation classes modulo $\mathbb Z d$. In each class write a phase
position as $b+a_i d$ and let $h=\max_i a_i$. The five required black rays
start at $b+h d$; their finite boundary assignments are exactly $A$.
A phase read in epoch $j<n$ has ray index $j+a_i-h<n$, so every nonnegative
index is present in the truncation and every negative index uses the same
boundary assignment. Literal coupling proves the entire prefix. At epoch
$n$, phase zero first reads the missing site $nd$, which is white.
The fixed class table and its infinite parity intervals are checked exactly.

Thus $D(u^n)=66n$, disproving a seed-independent constant for window eight
and every smaller window. For $k\ge9$, a $k$-legal word of length at least
$k$ is nine-legal; shorter words have deficit at most $29(k-1)$.
Therefore a finite uniform constant exists exactly when $k\ge9$.
This threshold concerns finite-word mass deficit, whereas eventual
attraction is governed by the ten-factor theorem.


\section{Exact first-violation
macros}\label{exact-first-violation-macros}

The preceding decision can be strengthened to return the first event
itself and its complete post-event state. This requires a finite-path
argument inside the nonsimple component, rather than an assumption of
periodicity there.

\textbf{Theorem F (exact event and finite prefixes).} A complete finite
field may be given as the XOR of \(M\) finite arithmetic segments whose
strides belong to
\[\{(0,0),(\pm2,\pm2),(\pm2,0),(0,\pm2),(\pm4,0),(0,\pm4)\}.\] Zero
stride denotes a singleton. The first future \(J_{10}\), or its
permanent absence, is decidable in polynomial bit time in this
representation. A finite answer returns its exact time and full state
just after the event. One call uses at most five periodic plans and 1111
literal main-orbit updates, adding at most 1530 segment terms.

Let \(W_0\) be the diameter in the maximum norm of the actual initial
support and head. Conditional on \(n\) events existing, their post-event
time satisfies \[s_n\le n(212W_0+854199)+415308n(n-1). \tag{MACRO1}\]
Every intermediate complete state lies in the initial convex hull
enlarged by \([-1959n,1959n]^2\). The first \(n\) events and their
complete encoded endpoint are computable in
\(\operatorname{poly}(M+n+b)\) bit operations, with \(n\) measured
numerically and \(b\) the original coordinate and length bit size.

These macros preserve finite histories exactly. Their lift counts fold
starts within one event-free traversal; a fresh event begins a fresh
traversal. The following proof supplies both the finite graph coverage
and the arithmetic continuation needed for this stronger conclusion.

\subsection{Pinned base and finite-path coverage}\label{macro-3}

The base contains every physically feasible 150-letter word all of whose
ten-factors are in L10. Its 6605 edge words generate a graph G with 6501
vertices, the 149-letter prefixes/suffixes. Physical feasibility is
exactly alternation of requested turns on consecutive reads of the same
physical cell. This uses the established forbidden-word method
here.

Every actual event-free continuation, once its 149-turn history is a G
vertex, follows G. At most 149 nonviolating ordinary updates are needed
to obtain such a vertex: afterwards the last 149 turns are wholly new,
and if the next turn is legal their feasible 150-word is in the complete
enumeration. If the next turn is illegal, it is returned as the event
instead. This covers input time zero without inventing a periodic
negative-time past.

The pinned graph has simple W and T components of lengths 104 and 24,
and one component H with 366 vertices and 370 internal edges. The prior
explicit quotient maps H to phases 0,\ldots,101. There is an edge
(j,B\_j)-\textgreater j+1 mod 102, plus (0,L)-\textgreater75, where B is
a 102-letter fold word beginning R. The short loop is A=L
B\_75\ldots B\_101, length 28. Their literal geometry is

\begin{verbatim}
A: displacement (0,-4), heading restored;
B: displacement (1,11), heading reversed.
\end{verbatim}

Every internal H edge is checked against this quotient, not merely
sampled. The already supplied all-parameter lemma is

\begin{verbatim}
B A^n B A^m B is physically infeasible for ALL n,m>=0. (MACRO2)
\end{verbatim}

Its explicit consecutive same-cell/same-R conflicts and their parameter
coverage are in Subsection \ref{ten-4}. A finite word need not extend to
an infinite legal word for (MACRO2) to apply. This is the essential
finite-path rather than infinite-cycle input.

The base edge set and quotient have frozen SHA256 values in the source
witness. The event runtime regenerates both and checks exact equality.

\subsection{A complete bounded-fold lift}\label{macro-4}

Outside H keep each vertex v with counter 0. Inside H use (v,c),
c=0,1,2,3, except that (v,3) with quotient phase 0 is forbidden. On an
internal H edge increase c exactly when its quotient starts the long
loop, namely phase 0 with appended R. All other internal edges keep
c.~On entry into H start c=0; on exit forget c.~Disallow c\textgreater3
and disallow arrival at phase 0 with c=3.

This permits two complete B loops and a third incomplete B loop; a third
*completed* loop is excluded. An initial partial loop is not charged, so
no assumption is made about the path before H entry. Between three
long-loop starts at phase zero there are powers of A, and completing the
third yields (MACRO2). Thus every real finite event-free path lifts. H
cannot be revisited after leaving it because G's SCC condensation is
acyclic. Counter reset on reentry cannot therefore mask additional
physical folds.

The lift may still contain impossible paths. It is an exclusion-safe
overapproximation, not a claim that every lifted path is physical.

Complete reconstruction yields 7595 vertices and 7763 edges. Its ONLY
cyclic SCCs are simple cycles of lengths

\begin{verbatim}
24, 104, 28, 28, 28.
\end{verbatim}

The 28-cycles are A phases at counters 0,1,2. The 24- and 104-cycles are
the prior T and standard words. No claim follows from cycle lengths
alone; the words and their nonzero displacements are checked too.

Assign to each condensation component the three weights (one
cyclic-component visit, one noncyclic-vertex update, period):
respectively (1,0,P) on a cyclic component, and (0,1,0) on a noncyclic
singleton. Backward longest-path values are nonnegative and, over every
condensation edge C-\textgreater D, satisfy

\begin{verbatim}
rank(C) >= weight(C)+rank(D)
\end{verbatim}

coordinatewise. Terminal ranks dominate their own weights. Internal
simple-cycle edges keep the same rank. All 7763 lifted edges are
checked. The maxima of the three ranks are

\begin{verbatim}
5, 956, 212.                                          (MACRO3)
\end{verbatim}

Python uses the base's iterative two-pass SCC routine; the independent
native checker reconstructs the lift and uses Tarjan SCCs. Sorted full
edge and rank hashes agree. The source witness fixes the complete edge and rank hashes. Integer inequalities, complete reconstruction and (MACRO2), not a
floating-point solver, carry the proof.

\subsection{Exact physical acceleration on a simple
component}\label{macro-5}

At a cyclic vertex propose its internal periodic word u, of length P and
nonzero displacement d, restoring heading. Let q\_i be its phase read
positions in the actual pose. At q\_i+kd the required old-field
contribution is determined by

\begin{verbatim}
actual required turn bit = 1_B(q_i+kd)
   XOR #{(j,s): q_j-q_i=s d,
                s>0 or (s=0 and j<i), 0<=s<=k} mod 2. (MACRO4)
\end{verbatim}

This computes the literal first mismatch by induction up to that
mismatch; it is not a finite-repeat heuristic. Prior-write changes occur
at finitely many epoch thresholds.

For a finite segment a+jv, 0\textless=j\textless L, its intersection
with a phase ray q+kd is empty, a singleton, or a finite arithmetic
interval in k. Nonparallel directions use an exact determinant. For
parallel directions solve the one-dimensional integer congruence after
dividing by the greatest common divisor. In the closed stride family its
epoch modulus is 1 or 2. This modulus matters: the 24-mode has axis
stride 2 and the 28-mode axis stride 4. Treating every parallel
intersection as a full interval would be wrong.

Split epochs into k=2z and k=2z+1. Each contribution is then an ordinary
finite interval in z; each write threshold is a final interval. Sort
parity-change endpoints separately for the two classes, including z=0
and the final infinite interval. This returns the exact earliest
mismatch, or certifies none. All arithmetic is on integers of polynomial
bit size.

Before a finite mismatch at t=nP+j, append P finite write segments for
the full periods (if n\textgreater0), and j singleton writes for the
partial period. These are XOR contributions; EVERY skipped write is
accounted for. There are at most 2P-1 new terms. Shift the history
vertex around its simple cycle by t mod P. Read the real cell and
execute its actual update. If J10 is true, return the complete
post-state. Otherwise the real edge must leave this simple component:
the unique internal edge has the different predicted label. Continue in
the lifted graph.

If no mismatch exists, the actual infinite continuation equals this
cycle word. A finite field cannot realize the 24- or 28-word forever:
their stabilized first-read classes require respectively two and five
infinite initial black rays, as proved in the prior theorem. Hence only
the standard cycle can give permanent absence. All earlier actual
updates were individually checked or were legal cycle edges, so the
certificate says no J10 at ANY nonnegative time, not only after its
standard boundary.

A component is never revisited after exit. By (MACRO3), one call
therefore uses at most five plans, at most 956 noncyclic updates, at
most five cycle exits, at most 149 warm-up updates, and one final event
update. A safe literal bound is 1111. The sum of visited cycle periods
is at most 212, hence

\begin{verbatim}
new terms <=149+956+2*212+1=1530.
\end{verbatim}

The 2P allocation includes a possible non-event cycle exit. No skipped
write or exit update is omitted. Periodic proposals and inverse-history
computations use a copy until the actual endpoint is determined.

\subsection{A generic finite periodic-plan halo lemma}\label{macro-6}

This lemma concerns ANY supplied periodic turn word of length P
restoring heading with d!=0, from a complete finite source Y=(B,p,h).
Suppose its first mismatch is at time kP+i, position e=p+q\_i+kd, with
phase coordinates relative to p.

Let L be the largest positive translation threshold s with q\_j-q\_i=s
d, or zero if none. Since a period has P unit moves,

\begin{verbatim}
||q_i||_infinity<=P,   L||d||_infinity<=P.
\end{verbatim}

If k\textless=L then
\textbar\textbar e-p\textbar\textbar\_infinity\textless=2P. If
k\textgreater=L+1, all planned previous-write parities at phases i in
epochs k-1 and k have stabilized. If the stable required old-field bit
is 0, mismatch means e is an original black cell. If it is 1, the
preceding occurrence matched, so e-d is an original black cell. Thus in
every case

\begin{verbatim}
e belongs to conv(B union {p})+[-2P,2P]^2.             (MACRO5)
\end{verbatim}

Here B is the FULL source field before the plan, not an unknown active
subset. For k\textgreater=1, any earlier phase point q\_j+ell*d differs
from (ell/k)(q\_i+kd) by at most 2P in infinity norm. For k=0, all reads
are within P of p.~Therefore the entire finite planned read path through
e is inside

\begin{verbatim}
conv(B union {p})+[-4P,4P]^2.                         (MACRO6)
\end{verbatim}

Its ensuing literal update enlarges this by at most one. Unchanged
source cells remain present outside the read footprint; (MACRO6) never
authorizes deleting them.

For radius R=max\_\{b in
B\}\textbar\textbar b-p\textbar\textbar\_infinity (zero for empty B),
(MACRO5)'s more precise alternatives give: either
k\textless=L\textless=P, or
k\textbar\textbar d\textbar\textbar\textless=R+\textbar\textbar d\textbar\textbar+P.
Therefore the first mismatch time is strictly less than

\begin{verbatim}
P(R+P+2).                                           (MACRO7)
\end{verbatim}

These are conditional-on-first-mismatch bounds. A valid infinite plan
has no finite halo bound.

\subsection{All-input J10 geometry, time and bit
complexity}\label{macro-7}

Sum (MACRO6) over all finite periodic portions of one event call. The
period sum is at most 212; all actual literal updates number at most
1111. Convex Minkowski enlargements compose, giving

\begin{verbatim}
H=4*212+1111=1959.
\end{verbatim}

Every intermediate complete state lies in the initial complete convex
hull enlarged by H; in particular its width is at most W0+2H and its
radius about its current head at most that width. Applying (MACRO7) to
every periodic part, using P\textless=104, and including all literal
updates, gives post-event duration at most

\begin{verbatim}
212(W0+2H+106)+1111 =212W0+854199.                   (MACRO8)
\end{verbatim}

After n-1 event calls width is at most W0+2H(n-1). Summing (MACRO8)
gives (MACRO1), since 212H=415308. No theorem bounds n.~In particular an
infinite chain is not classified by the finiteness of each call.

The support of every newly appended segment is made of actual read
positions. Thus all its endpoints lie in the same enlarged initial hull,
even if it later cancels other terms. After n events, all coordinate and
segment-length bit sizes are O(b+log(n+1)); event times also have
O(b+log(n+1)) bits by (MACRO1). The number of terms is at most M+1530n.
Each call performs a fixed number of ray/segment intersections, endpoint
sorts and fixed-length literal or inverse operations. Exact set/list
arithmetic therefore proves poly(M+n+b) bit complexity. Fixed constants
may be large; this is not a benchmark or a claim polynomial in log n.

Together with the already proved H(X) iff finitely many J10 events,
(MACRO1) implies that every hypothetical finite-support HC
counterexample, of EITHER the nonstandard-periodic or nonperiodic type,
has at least n intrinsic J10 events before every time T with (MACRO1)'s
right side \textless=T. In particular its liminf event count divided by
sqrt(T) is at least 1/sqrt(415308). This counts intrinsic events;
interval-contained ten-factor counts differ only in the first nine
boundary positions.

This is a lower rate requirement, not an upper event budget, a
NO-HIGHWAY certificate, or a contradiction. The mass-defect lower bound
and this timing bound can coexist with infinite renewal.


\section{Geometry of finite legal intervals}\label{sec:silent-geometry}
The local language controls the geometry of an entire legal interval, even
when its duration is unbounded. The coordinate selected in the next theorem
is fixed for the whole interval.

\textbf{Theorem (endpoint strip).}
Let $p_0,\ldots,p_L$ be the head positions of a physically feasible finite
word whose contained ten-factors are standard. There is an $i\in\{x,y\}$ such that
\[
 \min(p_{0,i},p_{L,i})-43\le p_{t,i}\le
 \max(p_{0,i},p_{L,i})+43\qquad(0\le t\le L).
\]
More precisely, after translating $p_0$ to zero, one of the four coordinate
normals $n$ satisfies $-43\le n\cdot p_t\le n\cdot p_L+43$ throughout.

\textit{Proof.}
Use the bounded-fold lift established for the event macros. Adjoining the
incoming heading gives a directed graph $G$ with 30,380 vertices and 31,052
edges. A vertex records the length-149 history, the fold count and the actual
heading; an edge appends its literal turn and carries its unit displacement.
Every feasible continuation after the first 149 turns lifts to this graph.
The all-parameter exclusion of $BA^nBA^mB$ supplies coverage of the fold count.

For a cyclic strongly connected component $C$, let
$M(C)=\{n:n\cdot d(C)\ge0\}$, where $n$ ranges over the four coordinate
normals and $d(C)$ is the displacement of its simple cycle. Give an acyclic
component all four normals. The finite certificate starts a mask at every
component and propagates every possible intersection along every condensation
edge. All 88,024 resulting component/mask states have nonempty masks.
Closure under propagation proves by induction that every finite path has a
common allowed normal. An independent work-list calculation on the original
graph agrees at all 90,256 vertex/mask states.

For each normal delete the cyclic components that omit it. On the resulting
$G_n$, the certificate gives an integer potential $H_n$ with
\[
 0\le H_n\le33,\qquad
 n\cdot\Delta(v,u)\ge H_n(u)-H_n(v)
 \quad\text{for every edge }v\longrightarrow u.
\]
There are 30,714 inequalities for each normal. Telescoping bounds every
subpath's retreat by 33. These inequalities and the nonempty-mask closure,
rather than the method used to find the potentials, are the finite premises.

Short words of length at most 149 have coordinate radius at most 24 by
complete feasible-word enumeration. For a longer word put $q=p_{149}$,
$v$ equal to its startup lift state, and $z_t=n\cdot p_t$. The startup
certificate checks, for every such state and every allowed normal,
\[
 \min_{t\le149}z_t\ge-43,\quad z_{149}-H_n(v)\ge-43,\quad
 \max_{t\le149}(z_t-z_{149}+H_n(v))\le43.
\]
For $t\ge149$, telescoping gives $z_t\ge z_{149}-H_n(v)\ge-43$
and $z_t\le z_L+33$. For $t\le149$, use the first startup inequality
and $z_L\ge z_{149}-H_n(v)$ in the third. Short words use the sign
with $n\cdot p_L\ge0$. This proves the oriented assertion and the strip.
\hfill$\square$

\textbf{Corollary (localized events).}
Let $K$ be a nonempty finite set with bounding rectangle
$[x_-,x_+]\times[y_-,y_+]$. Between any consecutive chain of intrinsic
$J_{10}$ events with read positions in $K$, all head positions lie in
\[
 \{x_- -44\le x\le x_+ +44\}\ \cup\
 \{y_- -44\le y\le y_+ +44\}.
\]
The conclusion includes the outgoing move of the last event and is independent
of the number of events in the chain. Indeed, each intervening silent word
starts one step from $K$ and ends in $K$; apply the theorem to that contiguous
word and include the event moves. If infinitely many events are eventually
localized in $K$, the entire late head trajectory lies in this same cross.

This geometry complements the complete-source bounds used by event prediction.
A legal interval occupies one endpoint strip while it may travel arbitrarily
far along it. The folded ribbon in Appendix~\ref{app:folded-ribbon} realizes precisely this distinction.


\section{Finite-horizon rigidity and silent supply}\label{sec:silent-supply}
The infinite-future attraction theorem does not apply to a legal fragment that
may later violate the language. We now bound the nonstandard part of every
finite ten-legal physical word, without assuming a legal continuation.
For such a word $w$, let $c(w)$ count distinct sites whose first requested read
is L. The alternating-read criterion makes this exactly the minimum initial
black mass realizing $w$. In particular $c(w)\le m$ for a segment starting
with a complete $m$-black field.

\textbf{Theorem (one contiguous standard core).}
Every finite ten-legal physical word has a decomposition
\[
 w=u(W_r)^k v,\qquad k\ge0,\qquad
 |u|+|v|\le114c(w)+154693,
\]
where $W_r$ is a cyclic phase of the standard word. If $\Delta$ is the total
head displacement, then
\begin{equation}\label{eq:silent-supply}
 |w|\le52\|\Delta\|_\infty+114c(w)+223437.
\end{equation}
The coefficient 52 is optimal; the other constants are sufficient bounds.
The core is one contiguous part of the actual trace, and it may be empty.

\subsection{Component order and transverse supply}
The finite-path fold cover established earlier represents every finite legal
word by at most five cyclic stages, with periods $24,104,28,28,28$.
The acyclic steps number at most 956 and the total cycle lengths at most 212.
Allowing the initial 149 letters, partial cycle remainders and exits leaves
at most
\[
 F=149+956+212+5=1322
\]
literal letters outside whole powers. The complete condensation graph checks
the additional order: at most one $T$ power precedes at most one $W$ power,
followed by at most three $A$ powers. A departed $W$ component cannot be
reentered within the same legal word. This order uses the all-parameter
three-fold obstruction, so it covers arbitrarily long finite paths.

Every phase of $T^k$ requires at least $2k$ distinct initially black read
sites, and every phase of $A^k$ requires at least $5k$. Each whole period
of either mode changes black mass by $+4$. These supply counts follow by
listing phases with no earlier same-cell predecessor in any translating
copy. If $q_i$ and $q_j$ are phase positions and $d$ the full drift,
an earlier copy can hit phase $i$ precisely when $q_j-q_i=s d$ with
$s>0$, or $s=0$ and $j<i$, subject to the copy being present.
The fixed phase tables check all rotations and starting phases.

A $W$ power reads an eleven-level diagonal strip, whereas an $A$ power
reads a four-level axis strip. Any translated pair intersects in at most
$11\cdot4=44$ cells: an axis coordinate and a diagonal coordinate determine
at most one lattice site. Thus an arbitrarily long standard block can supply
only 44 sites directly to each later axial block.

\subsection{Complete-field cofactor accounting}
Let $k_T,k_W,k_1,k_2,k_3$ be the whole-period counts, inserting zeros for
absent stages, and put $q=m+F$. Initially $2k_T\le q$.
Let $D$ be the complete net toggle set of the $W$ block. After that block
write the actual field as $B=R\mathbin{\triangle}D$. The cofactor $R$ is
an algebraic device, not a substituted physical state. Every cell of $D$
continues to participate in the actual evolution. For any toggle set $Z$,
\[
 |R\triangle Z|-|R|
 \le |(R\triangle D)\triangle Z|-|R\triangle D|+2|Z\cap D|.
\]
This holds pointwise: outside $Z\cap D$ the mass changes agree, and at an
intersection site they differ by at most two. Consequently an $A^k$ block
increases cofactor mass by at most $4k+88$, and at most 44 of its required
first-L sites can be supplied by $D$. Allocating all $F$ literal flips from
the start gives
\[
 5k_i\le3q+4\sum_{j<i}k_j+88(i-1)+44,\qquad i=1,2,3.
\]
Iteration yields
\[
 \sum_{i=1}^3 k_i\le\frac{453}{125}q+\frac{15004}{125}.
\]
The number $E$ of letters outside the whole standard power therefore obeys
\[
 E\le F+24k_T+28\sum_i k_i
 \le\frac{14184}{125}m+\frac{3867322}{25}
 \le114m+154693.
\]
Applying this bound to the minimal first-read source proves the core claim.

The common-normal certificate from the endpoint-strip argument supplies one
coordinate normal $n$ whose projection on every cyclic drift of this path
is nonnegative. If the standard power occurs, its projection is 2 per
period, while all literal letters together contribute at least $-F$.
Hence $2k_W\le\|\Delta\|_\infty+F$. Combining this with the bound on $E$
proves~\eqref{eq:silent-supply}, since $52F=68744$. Without a standard
power, the bound on $E$ already suffices. Standard prefixes $W^k$ from a
fixed finite seed have length $104k$, displacement norm $2k$ and bounded
first-black cost; they rule out any smaller displacement coefficient.

\subsection{Closed returns, lifetimes and a one-jump algorithm}
A closed legal return from a complete $m$-black field has length at most
$114m+223437$. If a cell is born black at time $b$ and next read at $e>b$,
and the word on $[b,e)$ is ten-legal, its endpoints coincide. Thus
\[
 e-b\le114|B_b|+223437.
\]
The mass belongs to the birth checkpoint. Repeated use of this inequality
does not replace changing checkpoint masses by the original input mass.

There is also an exact first-$J_{10}$ algorithm using one standard-periodic
query and at most $114m+154694$ ordinary forward updates on an explicit
complete $m$-point source. Recover the initial nine turns by inverse dynamics
on a copy of the full field. Execute literal updates until the first event,
or until the actual 149-letter history enters the standard component.
At that point the complete-field ray-parity query either certifies a
permanent standard future or gives its first mismatch. Retain every skipped
write as finite XOR segments and continue literally until the first event.
The component order forbids a second standard stage before that event.
The core bound limits all remaining literal work and proves totality.

The returned endpoint contains the original points, $O(m+1)$ point changes
and at most 104 finite translating segments. Coordinates and elapsed time
may be arbitrarily large binary integers; the algorithm has polynomial bit
complexity in the explicit sparse input. This does not improve the older
five-plan algorithm's guarantee for compressed inputs representing huge
black mass. Both algorithms retain all unread debris.

For a fixed finite target $K$, put
$R=\max_{q\in K}\|q-p_0\|_\infty$. If the first future event lies in $K$,
its pre-update index is at most $52R+114m+223437$. This bounds a first-event
query, not later returns to $K$. No total event budget or cross-event
descending rank follows.


\section{Limitations}
The attraction theorem assumes a legal future; the mass inequalities apply
to arbitrary finite trajectory intervals. Together they give local-event
algorithms and necessary density laws. They provide no upper bound on the
total number of events of an arbitrary finite seed. The ten-factor threshold
for attraction is sharp over arbitrary initial colourings; its optimality
on finite black supports is a separate question. The mass theorem has the
distinct sharp window threshold nine.

The coefficient twelve in the attraction bound and the constant 322 in the
ten-factor deficit bound have matching witnesses. The attraction intercept
212 and the penalty coefficient 351 are sufficient constants whose optimality
is open. Algorithmic costs distinguish explicit point sets from compressed
XOR segments and depend polynomially on the numerical event count, rather
than its bit length. Lifetime bounds use the complete mass at the named
checkpoint; repeated application does not replace it by the initial mass.
The memory families have finite prescribed prefixes and recoveries. Whether
one finite seed can sustain autonomous infinite renewal remains a global
dynamical question.

\section{Reproducibility and
conclusion}\label{reproducibility-and-conclusion}

The source package provides two complete finite kernels. The attraction
kernel regenerates the 6,605 physical length-150 words, the quotient of
the nonsimple component, and the exact path and supply bounds. The mass
kernel regenerates all 29,559 length-300 edges, checks every integer
potential inequality, and replays the equality witness. Python and C++
implementations are compared, and damaged certificates are rejected. The
event lift additionally checks all 7763 edges and component-rank maxima
5, 956 and 212. The all-parameter foldback exclusion, seed induction and spiral separation
are proved in the text; finite geometric checks verify their stated
numerical premises.

The attraction edge list and quotient map are
\path{verification/ten/certificates/windows10.txt} and
\texttt{quotient.tsv}; the mass witness is
\path{verification/mass/witness.json}. These are paths within the public
source package pinned below. For the mass graph, Python uses breadth-first
last-read maps and C++ depth-first reversible colour assignments; their sorted
edge and potential hashes agree. The event lift is reconstructed using two-pass
and Tarjan SCC algorithms, with identical sorted edge and rank hashes.

The memory constructions additionally reconstruct every affine read interval,
compare complete endpoint fields, verify the all-future standard recovery,
and check every alignment in the reset lower bound. Their universal parameter
coverage comes from the finite arrangement proof, with separate literal
simulations and damaged-witness controls.

The silent-supply checker independently rebuilds the 6,605-word cover,
component order and common normal, checks every phase supply and transverse
intersection bound, and compares 174 complete first-event endpoints.
The nine-factor checker independently enumerates all 164,283 length-150
words, checks 9,550 product-edge inequalities and the complete condensation
budget, and verifies the all-parameter eight-factor truncation.

Run \texttt{python3\ factor\_language/reproduce.py} from the source
root. Each checker implements the finite coverage or integer inequality
stated in its corresponding proof, and records its source identifiers and
verification results. The analytic arguments establish coverage for all
parameters; the exact computations establish their finite premises.

The computational sources and certificates are fixed at
\href{https://github.com/kehao95/langtons-ant-core/tree/e93252230ea9109908c23a009824d72d30fb91e2/factor_language}{repository commit \texttt{e932522}, factor project}.
The \texttt{sources.json} manifest lists their hashes. The accompanying
source archive includes those finite kernels and this article's source.

The two proof routes expose complementary structure: a legal infinite
future has rigid eventual motion, while an arbitrary finite interval
must pay for its mass deficit with local violations. Their combination
connects a local signature of P104 with both a decidable event and a
quantitative all-input law.

\appendix

\section{Memory produced during legal intervals}
\label{sec:silent-memory}
A finite ten-legal interval can perform substantial computation on its
complete input field. We give an explicit family that manufactures black
lifetimes later consumed at violations, followed by a related family with
an exact autonomous recovery and a quantitative reset cost. These
constructions identify the spatial information used by the event macro.

\subsection{A legal producer and a finite consumer}
Use the words \(A,B\) of the two-loop component and set
\[
 V_n=B[9:]A^nBA^{n-1},\qquad
 P=\mathrm{RLRRRRLRRRRLLLLRLRRL},\qquad n\ge1.
\]
Here \(B[9:]\) removes its first nine letters. For finite templates
\(L_0,M_0,H_0\), define
\[
 B_n=L_0\cup\bigcup_{j=0}^{n-1}(M_0+(0,4j))
              \cup(H_0+(0,4n)).
\]
The complete templates are
\begin{verbatim}
L0 = {(-1,0),(-1,3),(-1,6),(-1,7),(0,-1),(0,3),(0,5),(0,7),
      (1,-1),(1,2),(1,3),(1,4),(1,6),(1,7),(2,-2),(2,2),
      (2,4),(2,5),(3,-2),(3,1),(3,2),(4,-1),(4,0)}
M0 = {(-1,10),(-1,11),(0,8),(0,9),(0,11),
      (1,8),(1,10),(1,11),(2,8),(2,9)}
H0 = {(-3,11),(-3,12),(-2,9),(-2,10),(-2,13),
      (-1,9),(-1,11),(-1,12),(-1,13),(0,8),(0,9),(0,12),
      (1,8),(1,10),(1,12),(2,8),(2,9),(2,12)}
\end{verbatim}
Their vertical ranges are disjoint. Start the ant at \((-1,0)\), heading
west, with precisely \(B_n\) black; its mass is \(10n+41\).

\textbf{Theorem (silent production and later consumption).}
The trajectory begins with \(V_nP^{n+1}\mathrm{RLL}\). Put
\(T_n=56n+167\). There is no intrinsic \(J_{10}\) at
\(0\le t\le T_n\); the first occurs at \(T_n+1\), at \((1,4)\), heading
east, and leaves the head at \((1,5)\), heading north. The largest visited
height before that event is \(4n+13\). Every initial black lifetime has
already been consumed before \(T_n\).

For \(0\le j\le n\), the following are consecutive white/black reads:
\[
\begin{array}{c|c|c|c}
\text{site}&\text{birth}&\text{consumption}&\text{age}\\\hline
(1,4+4j)&44+28j&56n+168+20j&56n+124-8j\\
(0,7+4j)&56n+146-28j&56n+186+20j&40+48j
\end{array}
\]
Every birth is in the \(J_{10}\)-free interval and every displayed
consumption is a \(J_{10}\). There are exactly \(9n+6\) violations during
\(P^{n+1}\). The first mismatch with the continuation \(P^\omega\) is at
relative time \(20n+22\) after \(T_n\).

\textit{Proof by a complete affine arrangement.}
Every phase of a variable block has a fixed horizontal coordinate, vertical
coordinate \(4an+b+4\epsilon k\), and repetition index
\(0\le k<cn+d\), where \(\epsilon=\pm1\). In a physical cell class
\((x,4r+s)\), its visits therefore occupy a half-open interval
\[
 \alpha n+\beta\le r<\gamma n+\delta.
\]
Their times are affine in \(n,r\). Within a repeated block the coefficients
are the same for every phase, so the constant term fixes chronological
order; distinct blocks retain their script order. Literal phases and input
points give singleton intervals. All 274 phase records are represented.

Collect the affine endpoints of source cells, visits, target cells and the
selected lifetimes. Their order changes only at pairwise rational crossings.
Splitting at the floor and successor of every crossing produces the complete
parameter partition \(n=1,2,3\) and \(n\ge4\). Within each slab, an affine
inequality is checked at its endpoints; the unbounded slab additionally
requires the appropriate slope sign. Between consecutive spatial endpoints,
the source bit and entire ordered read list are constant. Starting at the
source bit and toggling after each read checks physical feasibility for every
cell and every parameter in that slab. Including all target intervals also
checks the complete endpoint field, including sites not read by a block.

The declarative witness records complete low/bulk/high endpoint profiles.
The verified endpoint data are
\[
\begin{array}{c|c|c|c}
\text{prefix}&\text{time}&\text{pose}&\text{black mass}\\\hline
V_n&56n+167&(0,4),\mathrm N&18n+58\\
V_nP^{n+1}&76n+187&(0,4n+8),\mathrm N&22n+62\\
V_nP^{n+1}\mathrm{RLL}&76n+190&(0,4n+9),\mathrm W&22n+61
\end{array}
\]
For every initially black band, the arrangement also verifies a read in
\(V_n\), proving exhaustion of source lifetimes. On each selected lifetime
band, the last producer read is R and the first consumer read is L; their
affine times give the table above. Endpoint-and-slope identities cover the
entire band, including all intervening integer cells. There are 607 nonempty
affine bands and 27 selected lifetime bands in the exact certificate.

Nine inverse updates stay below height eight and recover the true prefix
\(B[:9]\). All ten-factors of \(A,B\) and their possible junctions occur
in \(BAABAA\); its 307 windows are legal. The final nine letters of \(A\)
and \(B\) agree, so the producer's outgoing context is independent of
\(n\). With that context, violations in the first \(P\) occur at indices
\(1,2,3,4,5,19\), and in later copies at
\(0,1,2,3,4,5,6,7,19\). This proves the first event and \(6+9n\) count.
The affine path endpoints give the exact maximum height. Finally the
following RLL differs from \(P\) at its third letter, giving the exact
mismatch time. \(\square\)

The construction supplies \(2(n+1)\) distinct black lifetimes produced
without a violation and later consumed at violations. At least \(n+1\)
have age at least \(48n+124\). The cell \((1,4)\) is born at the fixed
time 44 and survives until the first violation, with age \(56n+124\).
For any fixed \(D\), those sufficiently old consumptions are simultaneously
\(J_{10}\) and \(S_D\) events. Here \(S_D\) denotes a black read
whose cell was not visited in the actual preceding \(D\) updates; this
is a different predicate from a forbidden ten-letter factor.

The fixed source and event poses accompany an unbounded excursion. Thus a
first-\(J_{10}\) macro requires information about its complete source
geometry; its two endpoint poses alone do not bound its path or duration.
The earlier complete-source convex-hull bound provides exactly that geometry.
The black mass of this family grows with \(n\).

\subsection{The consumer's finite fuel budget}
Each copy of \(P\) restores the heading and moves by \((0,4)\). Its entire
read footprint has local heights zero through three, so consecutive copies
have disjoint read footprints. The first-L cells are \((1,0)\) and
\((0,3)\), and both must already be black at the beginning of the train.
Consequently a complete field of mass \(m\) supports at most
\(\lfloor m/2\rfloor\) consecutive whole copies of \(P\).

One copy leaves four net new black cells, but these lie behind the next
copy's required first reads. The exact distinction between total mass and
remaining fuel explains how the constructed finite consumer can increase
black mass while exhausting the input it needs to continue.

\subsection{Autonomous recovery of a folded ribbon}
A second family gives a full infinite continuation, complementing the
finite producer--consumer construction. Define
\[
\begin{split}
 C&=\mathrm{RLRLLRLLRLRLLRRRRLLRLLRRRRLLRRRRLLRLRRRRLRLL},\\
 D_0&=\mathrm{LLRRRRLRRLRRRRLLLLRLRRRRLRL},\\
 w_n&=CA^{n+6}BA^{n+2}D_0,\qquad |w_n|=56n+397,
 \quad n\ge0.
\end{split}
\]
Run the geometric plan from the origin facing north, and place initial
black cells precisely at sites whose first requested letter is L. All
other cells are white. This specifies a complete seed \(X_n\) with
\(10n+76\) black cells.

\textbf{Theorem (recovery and reset cost).}
The prescribed prefix is physically feasible. It writes the origin black
at time zero and next reads it at time \(56n+396\). Its intrinsic
\(J_{10}\) events occur at times zero through eight, followed by none
until time \(56n+397\). Its first permanent standard time and complete
future violation count are
\[
 \tau(X_n)=152n+2527,\qquad
 \#\{t\ge0:J_{10}(F^tX_n)\}=56n+926.
\]
Let \(Y_n=F^{56n+398}X_n\). Any literal repair of length \(\ell\) from
\(Y_n\) enabling the prefix \(w_m\), at any translated and quarter-turned
pose and with arbitrary unread debris, satisfies
\[
 \ell\ge\max(0,4m-136).
\]

\textit{Proof.}
The exact recovery is six alternating scans,
\[
 E_0 V_1^{n+5}E_1V_1^{n+5}E_2
 V_2^{n+4}E_3V_2^{n+4}E_4V_3^{n+3}E_5V_3^{n+4}E_6,
\]
where
\[
\begin{gathered}
 V_1=\mathrm{RLRLLRLR},\quad
 V_2=\mathrm{RLRLLRLLLLRRRRLR},
\\
 V_3=\mathrm{RLRLRLLLLRLRRRRLLLLRLRRR}.
\end{gathered}
\]
The seven fixed words \(E_i\), of lengths
\(3,258,2,142,449,566,542\), are the literal finite data in the supplied
\texttt{silent\_fold\_reset/grammar.json}. Their geometric phases and all
variable scans are checked from their letters. This recovery has length
\(96n+2338\), ends at \((8,7)\), east, with mass \(18n+301\), and preserves
every historical cell.

For these words, writing a cell as \((x,r-4k)\) gives affine visit intervals
whose endpoint slopes are zero or one. All pairwise endpoint crossings
occur by parameter one. Consequently the complete order cases \(n=0,1,2\)
cover every \(n\ge0\); equality types persist beyond the final crossing.
The 2,287 derived phase records check all requested colours in chronological
order and give exact complete fields. The full reduction and recovered
cap coordinates are included in the supplemental proof and mask certificate.

The final 208 letters are two copies of one standard phase, preceded by a
letter disagreeing with its backward extension. At that endpoint, all
black cells have horizontal coordinates between \(-14\) and 10. Each
future standard ray of drift \((2,2)\) therefore has only finitely many
possible source hits. Checking the source hits, self-write overlap thresholds,
and final infinite interval certifies the entire tail. The order-case
argument covers all their colours for every \(n\). Subtracting 208 from
the endpoint time gives the stated first entry. Each added pair of scans
contributes \(2(8+9+11)=56\) intrinsic violations; the exact zero-parameter
count is 926, including the initial inverse-history windows. The last
violation is at \(\tau(X_n)+7\).

For the reset bound, the source requirements \(I\) and output bulk \(O\)
on a four-row tile are
\[
\begin{array}{c|c|c}
x&I&O\\\hline
-5&00..&1100\\-4&0000&1111\\-3&0011&1100\\-2&1101&0001\\
-1&1011&1000\\0&1100&0011\\1&0000&1111\\2&..00&0011
\end{array}
\]
A dot imposes no requirement. The exact overlap table checks every vertical
residue, horizontal shift and parallel orientation, giving at least four
mismatches per complete target tile. Two boundary tiles can lose at most
eight from this lower bound; the 128 output cap cells can repair at most
128 further mismatches. This gives \(4m-136\). At perpendicular orientations,
the width-eight source meets at most three target tiles; the remaining
tiles each require ten black cells against white, giving
\(10\max(0,m-3)\ge4m-136\). Every repair update toggles one site, so it
needs at least as many updates as distinct necessary corrections.
\(\square\)

For a clean reset to exactly the complete seed of size \(10m+76\), the
output mass change from \(Y_n\) is \(10m-18n-52\). The mass--defect
inequality proved next consequently gives
\[
 351N_{10}\ge3\ell+468n-260m+1030.
\]
For \(m=n\ge34\), the flip lower bound yields
\(N_{10}\ge\lceil(220n+622)/351\rceil\). This last event bound concerns
a clean reset; the flip bound itself permits arbitrary unread debris.
Neither statement assumes that such a reset exists.


\section{Exact words}\label{exact-words}

The machine-readable words are in
\texttt{verification/ten/certificates/witness.json}; line breaks below
are not letters. All geometric statements use north at the origin unless
a phase placement is explicit.

\subsubsection{W (length 104)}\label{w-length-104}

\begin{Shaded}
\begin{Highlighting}[]
\NormalTok{RRLLLLRLRRRRLRRRRLLLLRLRRRRLRLLRRLLLLRRLLRRRRLLRRLRL}
\NormalTok{LRLLRLRLLRRRRLLRLLRRRRLLRRRRLLRLRRRRLRLLLLRRRRLRRLRR}
\end{Highlighting}
\end{Shaded}

\subsubsection{U (length 28)}\label{u-length-28}

\begin{Shaded}
\begin{Highlighting}[]
\NormalTok{RRLRLLLLRRRRLRRLRRRRLLLLRLRR}
\end{Highlighting}
\end{Shaded}

\subsubsection{T (length 24)}\label{t-length-24}

\begin{Shaded}
\begin{Highlighting}[]
\NormalTok{RRRRLLRLLRRRRLLRRRRLLRLL}
\end{Highlighting}
\end{Shaded}

\subsubsection{A (length 28)}\label{a-length-28}

\begin{Shaded}
\begin{Highlighting}[]
\NormalTok{LLRRRRLRRLRRRRLLLLRLRRRRLRLL}
\end{Highlighting}
\end{Shaded}

\subsubsection{B (length 102)}\label{b-length-102}

\begin{Shaded}
\begin{Highlighting}[]
\NormalTok{RRLLLLRRLLRRRRLLRRLRLLRLLRLRLLRRRRLLRLRRRRLRLLLLRRRR}
\NormalTok{LRRLRRRRLLLLRLRRRRLRLLLLRRRRLRRLRRRRLLLLRLRRRRLRLL}
\end{Highlighting}
\end{Shaded}

\subsubsection{C (length 89)}\label{c-length-89}

\begin{Shaded}
\begin{Highlighting}[]
\NormalTok{RRLRRRRLLLLRLRRRRLRRLRRRRLLLLRLRRRRLRLLLLRRRRLRRLRRR}
\NormalTok{RLLLLRLRRRRLRLLLLRRRRLRRLRRRRLLLLRLRR}
\end{Highlighting}
\end{Shaded}

\subsubsection{Z (length 104)}\label{z-length-104}

\begin{Shaded}
\begin{Highlighting}[]
\NormalTok{RRRRLLRLRRRRLRLLLLRRRRLRRLRRRRLLLLRLRRRRLRRRRLLLLRLR}
\NormalTok{RRRLRLLRRLLLLRRLLRRRRLLRRLRLLRLLRLRLLRRRRLLRLLRRRRLL}
\end{Highlighting}
\end{Shaded}

\clearpage
\section{First-black supply table for the
24-mode}\label{first-black-supply-table-for-the-24-mode}

The two rays are each advanced by 24. The finite times and the rays are
disjoint as read events and correspond to distinct first-read cells.

\begin{longtable}[]{@{}rllr@{}}
\toprule\noalign{}
Phase & Finite first-L times & Ray starts & Maximum \(L-12b_j(L)\) \\
\midrule\noalign{}
\endhead
\bottomrule\noalign{}
\endlastfoot
0 & 5,7 & 8,23 & 5 \\
1 & 3,4,6 & 7,22 & 3 \\
2 & 2,3,5 & 6,21 & 2 \\
3 & 1,2,4 & 5,20 & 1 \\
4 & 0,1,3 & 4,19 & 0 \\
5 & 0,2 & 3,18 & 0 \\
6 & 1 & 2,17 & 1 \\
7 & 0,7 & 1,16 & 0 \\
8 & 6 & 0,15 & 0 \\
9 & 5 & 14,23 & 5 \\
10 & 3,4 & 13,22 & 3 \\
11 & 2,3,11 & 12,21 & 2 \\
12 & 1,2,10,19 & 11,20 & 1 \\
13 & 0,1,7,9,18 & 10,19 & 0 \\
14 & 0,6,8,17 & 9,18 & 0 \\
15 & 5,7,16 & 8,17 & 5 \\
16 & 3,4,6,15 & 7,16 & 3 \\
17 & 2,3,5,14 & 6,15 & 2 \\
18 & 1,2,4,13 & 5,14 & 1 \\
19 & 0,1,3,12 & 4,13 & 0 \\
20 & 0,2,11 & 3,12 & 0 \\
21 & 1,10 & 2,11 & 1 \\
22 & 0,7,9 & 1,10 & 0 \\
23 & 6,8 & 0,9 & 0 \\
\end{longtable}

\section{Forbidden words for the nine-factor certificate}\label{app:nine-forbidden}
The fourteen exclusions used in Section~\ref{sec:nine-mass} are the seven
words below and their letterwise complements (interchange R and L).
For each word the table gives consecutive visits to the same cell demanding
the same turn. Indices start at zero; the geometric walk starts at the origin
facing north. This is a direct finite-word obstruction, valid on every
initial colouring. The listed conflicts have been checked with both integer
coordinates and complex-coordinate rotations.

\begin{center}\small
\begin{tabular}{rrrrrc}
\toprule Word & Length & Earlier & Later & Cell & Turn\\\midrule
1 & 161 & 0 & 160 & $(0,0)$ & R\\
2 & 161 & 0 & 160 & $(0,0)$ & L\\
3 & 153 & 0 & 152 & $(0,0)$ & R\\
4 & 153 & 0 & 152 & $(0,0)$ & L\\
5 & 253 & 0 & 252 & $(0,0)$ & L\\
6 & 253 & 0 & 252 & $(0,0)$ & R\\
7 & 181 & 0 & 180 & $(0,0)$ & R\\
8 & 181 & 0 & 180 & $(0,0)$ & L\\
9 & 217 & 0 & 216 & $(0,0)$ & R\\
10 & 217 & 0 & 216 & $(0,0)$ & L\\
11 & 309 & 0 & 308 & $(0,0)$ & L\\
12 & 309 & 0 & 308 & $(0,0)$ & R\\
13 & 237 & 0 & 236 & $(0,0)$ & R\\
14 & 237 & 0 & 236 & $(0,0)$ & L\\
\bottomrule\end{tabular}\end{center}

\noindent\begin{minipage}{\linewidth}
\noindent Word 1 (word 2 is its letterwise complement):
\begin{verbatim}
RRLRLLRLLRLRLLRRRRLLRLRRRRLRRLRRRRLLLLRLRRRRLRLLLLRRRRLRRLRRRRLLLLRLRRRR
LRLLRRLLLLRRLLRRRRLLRRLRLLRLLRLRLLRRRRLLRLRRRRLRLLLLRRRRLRRLRRRRLLLLRLRR
RRLRLLLLRRRRLRRLR
\end{verbatim}
\end{minipage}\par\medskip

\noindent\begin{minipage}{\linewidth}
\noindent Word 3 (word 4 is its letterwise complement):
\begin{verbatim}
RRLLRLRRRRLRLLLLRRRRLRRLRRRRLLLLRLRRRRLRLLLLRRRRLRRLRRRRLLLLRLRRRRLRLLRR
LLLLRRLLRRRRLLRRLRLLRLLRLRLLRRRRLLRLRRRRLRLLLLRRRRLRRLRRRRLLLLRLRRRRLRLL
LLRRRRLRR
\end{verbatim}
\end{minipage}\par\medskip

\noindent\begin{minipage}{\linewidth}
\noindent Word 5 (word 6 is its letterwise complement):
\begin{verbatim}
LRLRRRRLRLLLLRRRRLRRLRRRRLLLLRLRRRRLRLLLLRRRRLRRLRRRRLLLLRLRRRRLRLLRRLLL
LRRLLRRRRLLRRLRLLRLLRLRLLRRRRLLRLRRRRLRLLLLRRRRLRRLRRRRLLLLRLRRRRLRLLLLR
RRRLRRLRRRRLLLLRLRRRRLRLLLLRRRRLRRLRRRRLLLLRLRRRRLRLLRRLLLLRRLLRRRRLLRRL
RLLRLLRLRLLRRRRLLRLRRRRLRLLLLRRRRLRRL
\end{verbatim}
\end{minipage}\par\medskip

\noindent\begin{minipage}{\linewidth}
\noindent Word 7 (word 8 is its letterwise complement):
\begin{verbatim}
RRLRRRRLLLLRLRRRRLRLLLLRRRRLRRLRRRRLLLLRLRRRRLRLLLLRRRRLRRLRRRRLLLLRLRRR
RLRLLRRLLLLRRLLRRRRLLRRLRLLRLLRLRLLRRRRLLRLRRRRLRLLLLRRRRLRRLRRRRLLLLRLR
RRRLRLLLLRRRRLRRLRRRRLLLLRLRRRRLRLLRR
\end{verbatim}
\end{minipage}\par\medskip

\noindent\begin{minipage}{\linewidth}
\noindent Word 9 (word 10 is its letterwise complement):
\begin{verbatim}
RRLRLLRLLRLRLLRRRRLLRLRRRRLRRLRRRRLLLLRLRRRRLRLLLLRRRRLRRLRRRRLLLLRLRRRR
LRLLLLRRRRLRRLRRRRLLLLRLRRRRLRLLRRLLLLRRLLRRRRLLRRLRLLRLLRLRLLRRRRLLRLRR
RRLRLLLLRRRRLRRLRRRRLLLLRLRRRRLRLLLLRRRRLRRLRRRRLLLLRLRRRRLRLLLLRRRRLRRL
R
\end{verbatim}
\end{minipage}\par\medskip

\noindent\begin{minipage}{\linewidth}
\noindent Word 11 (word 12 is its letterwise complement):
\begin{verbatim}
LRRLRRRRLLLLRLRRRRLRLLRRLLLLRRLLRRRRLLRRLRLLRLLRLRLLRRRRLLRLRRRRLRLLLLRR
RRLRRLRRRRLLLLRLRRRRLRLLLLRRRRLRRLRRRRLLLLRLRRRRLRLLLLRRRRLRRLRRRRLLLLRL
RRRRLRLLLLRRRRLRRLRRRRLLLLRLRRRRLRLLRRLLLLRRLLRRRRLLRRLRLLRLLRLRLLRRRRLL
RLRRRRLRLLLLRRRRLRRLRRRRLLLLRLRRRRLRLLLLRRRRLRRLRRRRLLLLRLRRRRLRLLLLRRRR
LRRLRRRRLLLLRLRRRRLRL
\end{verbatim}
\end{minipage}\par\medskip

\noindent\begin{minipage}{\linewidth}
\noindent Word 13 (word 14 is its letterwise complement):
\begin{verbatim}
RRLRRRRLLLLRLRRRRLRLLLLRRRRLRRLRRRRLLLLRLRRRRLRLLLLRRRRLRRLRRRRLLLLRLRRR
RLRLLLLRRRRLRRLRRRRLLLLRLRRRRLRLLRRLLLLRRLLRRRRLLRRLRLLRLLRLRLLRRRRLLRLR
RRRLRLLLLRRRRLRRLRRRRLLLLRLRRRRLRLLLLRRRRLRRLRRRRLLLLRLRRRRLRLLLLRRRRLRR
LRRRRLLLLRLRRRRLRLLRR
\end{verbatim}
\end{minipage}\par\medskip


\section{References}\label{references}

\begingroup\small
\begin{enumerate}
\def\labelenumi{\arabic{enumi}.}
\tightlist
\item
  Anahí Gajardo. \emph{A symbolic projection of Langton's Ant (Extended
  Abstract)}. DMTCS Proceedings AB, 57--68 (2003).
  \href{https://doi.org/10.46298/dmtcs.2312}{doi:10.46298/dmtcs.2312}.
\item
  Anahí Gajardo, Victor H. Lutfalla and Michaël Rao. \emph{Ants on the
  highway}. Natural Computing 24, 497--509 (2025).
  \href{https://doi.org/10.1007/s11047-025-10018-9}{doi:10.1007/s11047-025-10018-9}.
  Author version: \href{https://arxiv.org/abs/2409.10124v1}{arXiv:2409.10124v1}.
\item
  Atharva Jillhewar. \emph{Finite-Support Periodic Highways of Langton's
  Ant: Necessary Conditions, Transverse Exclusions, and Exact Search}
  (2026).
  \href{https://doi.org/10.21203/rs.3.rs-10516269/v1}{doi:10.21203/rs.3.rs-10516269/v1}.
  Theorem 3.1 and Remark 3.2 supply full-drift visit chronology and seed
  reconstruction.
\item
  Anahí Gajardo and Jacques Mazoyer. \emph{One Head Machines from a
  symbolic approach}. Theoretical Computer Science 370, 34--47 (2007).
  \newline\href{https://doi.org/10.1016/j.tcs.2006.10.004}{doi:10.1016/j.tcs.2006.10.004}.
  Author manuscript (2005), Section 6, Proposition 8 and Theorems 1--2:
  \href{https://www.ing-mat.udec.cl/investigacion/paginas-pp/download.php?file=pp05-10.pdf}{Universidad de Concepción preprint pp05-10}.
\end{enumerate}
\endgroup

\end{document}
